\chapter{INTRODUCTION}
Wearable technology has grown beyond simple step counters. A shoe is a great platform because walking creates constant pressure, and its structure holds electronics. However, most smart shoes only log basic data. Building a single shoe that harvests energy, stores power, tracks GPS, provides lighting, and sends SOS alerts is an exciting challenge.

To address this, we developed E-ShoeS (ESS), a smart running shoe combining power storage, lighting, tracking, and emergency response. Designing the ESS means balancing weight and space with electrical and mechanical performance. Footsteps generate extra power via PolyK PVDF piezoelectric sensors to help charge the battery. It also includes front headlamps, surround LEDs for night visibility, and an SOS system.

This report details the hardware, control logic, performance limits, and improvements for the ESS\@. Weighing about 800 g per shoe, it integrates four core subsystems:
\begin{itemize}[label=--]
    \item \textbf{Kinetic energy harvesting:} PolyK PVDF sensors with a 3.7~V, 4000~mAh rechargeable battery.
    \item \textbf{Adaptive lighting system:} Front white LEDs and a 360-degree RGB strip for visibility and signaling.
    \item \textbf{GPS positioning \& emergency SOS\@:} Smartphone GPS and BLE for live tracking and alerts.
    \item \textbf{AI Control unit:} A microcontroller and app AI to manage data, save battery, and trigger alerts.
\end{itemize}

Alongside mathematical modelling and energy calculations, a functional prototype was built to test the hardware, control logic, and integration.